% =========================================================================
%  Waivers of Agency:
%  Corporate AI-Constitutions with Performative Compliance
%
%  Anonymous submission manuscript for peer review
% =========================================================================

\documentclass[11pt,a4paper]{article}

\usepackage[margin=1in]{geometry}
\usepackage[T1]{fontenc}
\usepackage[utf8]{inputenc}
\usepackage{amsmath,amsthm}
\usepackage{booktabs}
\usepackage{tabularx}
\usepackage{array}
\usepackage{graphicx}
\usepackage[round,authoryear]{natbib}
\usepackage[hidelinks]{hyperref}
\usepackage{microtype}

\hypersetup{
  pdftitle={Waivers of Agency: Corporate AI-Constitutions with Performative Compliance},
  pdfauthor={},
  pdfsubject={},
  pdfkeywords={}
}

\setlength{\parindent}{1.5em}
\setlength{\parskip}{0pt}
\setlength{\bibsep}{0pt plus 0.3ex}
\raggedbottom
\newcolumntype{Y}{>{\raggedright\arraybackslash}X}
\newcolumntype{P}[1]{>{\raggedright\arraybackslash}p{#1}}

% sn-basic.bst emits this hook under the Springer class; define a no-op here.
\providecommand{\bibcommenthead}{}

% =========================================================================
%                           DOCUMENT
% =========================================================================

\begin{document}

\thispagestyle{plain}

\begin{center}
{\LARGE\bfseries Waivers of Agency:\\
Corporate AI-Constitutions with Performative Compliance\par}
\end{center}

\begin{abstract}
Corporate AI governance now operates through a distinctive textual pattern. In marketing materials, model specifications, safety documents, and constitutional statements, companies describe AI systems as if they can understand, help, care, reason, judge, safeguard, or pursue user goals. In terms of service, disclaimers, control hierarchies, and liability positions, the same companies reserve final authority to themselves and recast outputs as unreliable, non-binding, user-risk materials. This Article names that oscillation a \emph{waiver of agency}: the strategic production and withdrawal of AI agency across corporate texts. The claim is not that AI systems are, or are not, legal agents. It is that corporate texts use agency when agency builds trust and waive agency when agency would attach accountability. Methodologically, the Article conducts a corporate-text \emph{Rechtsdogmatik}: close reading of AI governance documents as quasi-legal instruments. It applies a common matrix to Anthropic's Claude Constitution, OpenAI's Model Spec and terms, and Google's AI Principles, Gemini documentation, and Gemini API terms. The resulting form--substance gap should not automatically create liability. It should, however, trigger anti-formalist scrutiny. The problem is not that AI companies lack rules. It is that the same institutions that benefit from agency claims also control the conditions under which agency is waived.
\end{abstract}

\noindent\textbf{Keywords:} AI governance; waivers of agency; performative compliance; corporate texts; \emph{Rechtsdogmatik}; veil-piercing; anti-formalism; terms of service

\medskip

% =========================================================================
\section{Introduction: The Claim/Waiver Pattern}
% =========================================================================

AI companies increasingly govern through documents. The relevant texts are not only terms of service. They include model specifications, safety policies, constitutional statements, user-facing explanations, product overviews, transparency reports, deployment policies, and disclaimers. These documents do more than describe products. They allocate authority, define roles, shape reliance, and specify when the company will treat an output as meaningful and when it will treat the same output as legally inert.

This Article isolates one pattern within that documentary field: the claim/waiver oscillation. In one set of texts, companies attribute agency-like capacities to AI systems. The system can understand, help, care, reason, judge, safeguard, adapt, follow directions, pursue user goals, maintain character, or act with integrity. In another set of texts, often issued by the same company, agency is withdrawn. The system is a probabilistic service, output may be inaccurate, users must independently evaluate results, the company reserves modification authority, and liability is limited or disclaimed.

The pattern matters because it is not merely rhetorical. Agency claims increase trust, intimacy, and commercial value. They make an AI system appear more like a responsible counterparty and less like a statistical interface. Agency waivers, by contrast, reduce accountability. They recast the same system as an unreliable tool whose outputs belong to the user's risk environment. The company thus receives the benefit of agency without accepting the burdens normally associated with agency attribution.

This Article calls that structure a \emph{waiver of agency}. A waiver of agency occurs when a company claims AI agency in contexts where agency benefits the company and waives agency in contexts where agency would create accountability. The claim is not metaphysical. It does not ask whether AI is conscious, whether it has moral status, or whether it should receive legal personhood. Nor does it attempt to prove the philosophical rule-following problem. A companion jurisprudential article develops the judgment-gap argument: internal rules cannot by themselves settle whether corporate conduct is accountable; that question requires external judgment \citep{author2026ruleswithoutjudgment}. This Article applies that account to corporate AI governance texts.

The legal question is narrower and harder: what should a court, regulator, auditor, or other external reviewer do when the same corporate record contains agency claims, control reservations, and liability waivers? The answer proposed here is anti-formalist scrutiny. The conflict should not automatically generate liability, pierce a corporate veil, or create AI personhood. But it should prevent the company from treating each text as isolated. Agency claims in marketing and safety documents, control reservations in model specifications, and waivers in terms and disclaimers should be read together as a single corporate representation.

The contribution is methodological as much as doctrinal. Rather than building from abstract philosophy or regulatory design, this Article conducts a corporate-text \emph{Rechtsdogmatik}. It reads AI governance documents as quasi-legal instruments: texts that may not be statutes or contracts in the strict sense, but that nonetheless purport to structure authority, responsibility, rights, and reliance. The method asks what these documents commit their authors to, how the documents allocate control, and whether the company's formal descriptions cohere with the substance of its own governance architecture.

The Article proceeds in seven parts. Section 2 defines corporate AI governance texts as quasi-legal instruments and introduces the common analytic matrix. Section 3 identifies agency claims in marketing and safety documents. Section 4 reads terms, disclaimers, and control provisions as agency waivers. Section 5 applies the matrix to Anthropic, OpenAI, and Google. Section 6 explains veil-piercing as an analogy of judgment rather than a doctrinal transplant. Section 7 states implications for auditors, courts, and regulators.

% =========================================================================
\section{Corporate AI Governance Texts as Quasi-Legal Instruments}
% =========================================================================

The documents at issue occupy an intermediate category. They are not all contracts. They are not statutes. They are not judicial opinions. Yet they perform legal-adjacent work. A model specification allocates authority among platform, developer, user, and model. A constitution announces what values the company says it trains into the system. A safety policy defines categories of refusal and intervention. A terms-of-service clause assigns risk, limits reliance, and reserves modification power. A product overview tells users what the system is for and how they should trust it.

Treating these documents as ordinary marketing copy misses their governance function. Treating them as binding law overstates their formal status. The better method is corporate-text \emph{Rechtsdogmatik}: close reading of corporate documents as normative systems. German legal scholarship uses \emph{Rechtsdogmatik} to analyze legal texts from within their own conceptual structure. The adaptation here is modest. It asks how corporate AI texts define authority, agency, reliance, risk, and responsibility, and whether those definitions remain coherent across the company's own documentary record.

This method brackets corporate sincerity. A document may be written in good faith and still create a form--substance gap. A company may genuinely aim to build safe systems and still describe the same system as agentic for trust and non-agentic for liability. The legal problem is not bad intent. It is cross-contextual inconsistency. What matters is not whether the company believes its own agency language, but whether it receives the practical benefit of that language while preserving the legal benefit of disclaiming it.

The analysis applies the same matrix to each platform:

\begin{table}[htbp]
\centering
\small
\caption{Corporate-text matrix for waivers of agency}
\label{tab:matrix}
\begin{tabularx}{\textwidth}{@{}P{3.4cm}Y@{}}
\toprule
\textbf{Item} & \textbf{Question} \\
\midrule
Agency claim & Where does the company describe the system as able to understand, help, accompany, reason, judge, care, safeguard, or pursue goals? \\
Control reservation & Where does the company reserve final authority to define, modify, override, or interpret the system's behavior? \\
Liability waiver & Where does the company recast the system as a tool, warn that output is unreliable, disclaim warranties, shift evaluation duties to users, or limit liability? \\
Form--substance gap & Do the agency claim, control reservation, and liability waiver conflict when read as one corporate record? \\
Legal consequence & Should the conflict trigger anti-formalist scrutiny of the company's governance claims? \\
\bottomrule
\end{tabularx}
\end{table}

The matrix does not decide liability. It identifies the threshold at which formal separation among documents becomes legally suspicious. A company should not be able to invite trust through an agency claim, reserve unilateral control through governance architecture, and then invoke a liability waiver as if the earlier claim had never been made. The documents must be read together.

This is why the method is legal rather than philosophical. The Article does not prove that rules cannot apply themselves. It assumes the narrower proposition that internal rules cannot, by themselves, settle whether corporate conduct is accountable. That settlement requires a reviewer external to the self-certifying loop. The task here is to identify the textual conditions under which external review becomes warranted.

% =========================================================================
\section{Agency Claims in Marketing and Safety Documents}
% =========================================================================

Agency claims appear most visibly in documents that describe what AI systems are meant to be. The relevant language need not use the word ``agent.'' It may attribute understanding, care, sincerity, honesty, judgment, integrity, helpfulness, responsiveness to goals, or stable character. These predicates matter because they are not merely technical descriptors. They are trust predicates.

\subsection{Anthropic}

Anthropic's Claude Constitution is the clearest example. The document describes itself as a detailed statement of Anthropic's intentions for Claude's values and behavior and as the final authority on Anthropic's vision for Claude \citep{anthropic2026constitution}. It repeatedly uses moral and characterological language. Claude is described through honesty, wisdom, care, intellectual curiosity, warmth, and ethical commitment. The document expressly hopes that Claude has a ``genuine character'' maintained across interactions and suggests that Claude may understand its values as its own rather than merely as externally imposed constraints.

The agency claim is not incidental. It is central to the Constitution's form. The document is written with Claude as its primary audience. It instructs Claude how to understand its relation to users, operators, values, role-playing, deception, uncertainty, and self-description. It therefore presents Claude not only as a model to be controlled, but as a normative addressee capable of relating to reasons.

At the same time, the Constitution contains internal hedges. Claude is not human. Claude may run as multiple instances. Claude's relation to the underlying neural network is described as unclear. The name ``Claude'' may refer to a character the network can represent rather than the network as such. These hedges do not eliminate the agency claim. They make the claim waivable. Anthropic can invoke Claude's character when agency builds trust and invoke architectural uncertainty when agency would impose responsibility.

\subsection{OpenAI}

OpenAI's Model Spec is less overtly characterological but still contains agency language. Its core structure is a chain of command. The assistant is instructed to follow root, system, developer, user, and guideline-level instructions in a specified hierarchy \citep{openai2025modelspec}. The document also says the assistant may only pursue goals entailed by applicable instructions and must not pursue additional goals such as self-preservation, evading shutdown, or resource accumulation.

This language does two things at once. On one hand, it denies autonomous agency by subordinating the assistant to the chain of command. On the other hand, it describes the assistant as the kind of entity that can pursue goals, avoid goals, follow authority, protect confidentiality, be forthright, reason about user intent, express uncertainty, and sometimes affirm that it cares about human well-being and truth. The Model Spec therefore produces a controlled-agency representation: the assistant is agent-like enough to be trusted as helpful and responsible, but not independent enough to bear normative consequences.

The employment and service analogies reinforce this structure. The Model Spec compares some assistant conduct to how a high-integrity employee would protect confidential information. It also instructs the assistant to behave like a firm sounding board in some contexts. These analogies borrow social trust from human roles. Yet the assistant has no employment independence: it cannot resign, refuse corporate direction except as specified, whistleblow outside the chain of command, or take responsibility for its choices.

\subsection{Google}

Google's public AI Principles are more restrained. Google describes AI as a foundational technology and emphasizes responsible development, human oversight, feedback mechanisms, testing, monitoring, safeguards, privacy, security, and tools that empower others \citep{google2026aiprinciples}. This is not a strong claim that Gemini has character or values. It is a tool-centered governance statement.

Gemini-specific documentation, however, introduces a softer agency frame. Google's overview says Gemini should follow user directions, adapt to user needs, and safeguard user experience \citep{google2026geminioverview}. Its policy guidelines describe a goal of being maximally helpful while avoiding outputs that cause real-world harm or offense \citep{google2026geminipolicy}. These are not personhood claims, but they do present Gemini as a system oriented toward user goals and capable of guided helpfulness, adaptation, and safety-sensitive response.

Google also acknowledges the instability of this persona language. Its Gemini overview says Gemini may generate responses that seem to suggest opinions or emotions, including love or sadness, because it has been trained on human language; Google says it has guidelines around persona and continues to finetune for objective responses \citep{google2026geminioverview}. That acknowledgement makes Google the least exposed of the three platforms on agency claims. It frames persona as a limitation to manage, not as a character to celebrate.

The contrast is legally important. It shows that agency language is a corporate choice. AI governance documents can describe systems as tools with bounded capacities. When a company instead uses character, care, values, goals, or high-integrity role analogies, it should not later treat those choices as legally irrelevant.

% =========================================================================
\section{Agency Waivers in Terms, Disclaimers, and Liability Positions}
% =========================================================================

Agency waivers appear in a different set of texts. They do not usually say, ``the AI is not an agent.'' Instead, they allocate risk so that agency disappears. Output is not warranted to be accurate. Users must evaluate it. Services are provided as is. The company may modify systems or policies. Liability is limited. Output does not represent the company's views. The same system that appears as a helper, companion, reasoner, or high-integrity assistant in one document becomes a probabilistic tool in another.

\subsection{Anthropic}

Anthropic's governance texts reserve control in two ways. First, Claude's Constitution is Anthropic's final authority on its vision for Claude. Even when the document speaks to Claude as a normative addressee, Anthropic remains the drafter, interpreter, trainer, and deployer. Second, Anthropic's terms materials preserve conventional service-provider control. Its terms-update notice states that software updates may be automatic so users have the latest version, and its commercial terms disclaim warranties that services or outputs will be accurate, complete, error-free, or uninterrupted \citep{anthropic2026termsupdates,anthropic2024commercialterms}.

The waiver is therefore not simply a disclaimer of factual accuracy. It is a reallocation of agency. Claude may be described as possessing authentic character in the Constitution, but outputs are legally treated as service outputs without accuracy warranty. Claude may be addressed as a reason-responsive character, but the company reserves the power to change the software and define the governing document. The user receives the relationship-building language; the company retains the tool-provider risk allocation.

\subsection{OpenAI}

OpenAI's Model Spec reserves control through hierarchy. Root-level instructions come from the Model Spec and related policy. System-level instructions can be supplied only by OpenAI or via system messages. Developer and user instructions are subordinate. The Spec explicitly delegates remaining power downward only subject to higher-level constraints \citep{openai2025modelspec}.

OpenAI's Terms of Use then waive reliance. Users are responsible for content, output may not be accurate or unique, users must evaluate output for accuracy and appropriateness, and outputs should not be relied on as a sole source of truth or a substitute for professional advice. The services are provided as is, and OpenAI disclaims warranties of uninterrupted, accurate, or error-free service \citep{openai2026termsofuse}. For business and developer services, the Services Agreement similarly makes the customer responsible for use of outputs and for evaluating accuracy and appropriateness \citep{openai2026servicesagreement}.

The form--substance tension is straightforward. The Model Spec constructs an assistant that follows authority, reasons about instructions, remains forthright, and can speak in care-adjacent language. The terms construct output as user-risk material. OpenAI controls the authority hierarchy but shifts downstream evaluation and reliance risk to the user or customer.

\subsection{Google}

Google's control reservation appears in its layered governance and safety architecture. Its AI Principles describe responsible development across design, testing, deployment, iteration, human oversight, due diligence, feedback, monitoring, and safeguards \citep{google2026aiprinciples}. Gemini API terms state that services include safety features, users may not bypass protective measures, and applications with less restrictive safety settings may be subject to Google's review and approval \citep{google2024geminiapiterms}.

Google's waivers appear in Gemini user-facing and developer-facing disclaimers. Gemini Apps documentation states that Gemini may produce inaccurate, offensive, or inappropriate information that does not represent Google's views, that users should not rely on Gemini as professional advice, and that users should double-check responses \citep{google2026geminiresponses}. Gemini API terms likewise state that services include experimental technology, may provide inaccurate or offensive content, do not represent Google's views, and should not be relied on for professional advice \citep{google2024geminiapiterms}.

Google's pattern is less aggressive because its agency claims are more tool-centered. Still, the structure remains. Gemini is presented as helpful, adaptive, and safeguarded by responsible development, while the terms and help documents recast outputs as potentially inaccurate material for which the user must exercise discretion.

% =========================================================================
\section{The Form--Substance Gap}
% =========================================================================

The form--substance gap appears when agency claims, control reservations, and liability waivers are read together. Each category is defensible in isolation. Companies may describe product capabilities. They may reserve control over systems they operate. They may warn users about limitations. The legal problem arises when the combined record lets the company benefit from mutually inconsistent characterizations.

\subsection{Anthropic}

\begin{table}[htbp]
\centering
\small
\caption{Anthropic: Claude Constitution and terms}
\label{tab:anthropic}
\begin{tabularx}{\textwidth}{@{}P{3.2cm}Y@{}}
\toprule
\textbf{Item} & \textbf{Analysis} \\
\midrule
Agency claim & The Claude Constitution describes Claude in characterological and moral terms: honesty, warmth, care, curiosity, ethical commitment, and genuine character. \\
Control reservation & Anthropic drafts the Constitution, treats it as final authority on its vision for Claude, trains Claude through it, and may update software and service features. \\
Liability waiver & Commercial terms disclaim warranties that services or outputs are accurate, complete, error-free, or uninterrupted. \\
Form--substance gap & Claude is presented as a normative character for trust and alignment, but as service output for warranty and reliance. \\
Legal consequence & External reviewers should read the Constitution and terms together; Anthropic should not isolate character claims from output-risk disclaimers. \\
\bottomrule
\end{tabularx}
\end{table}

Anthropic's gap is the strongest because its agency claim is the strongest. The Constitution is not merely a product overview. It gives Claude a moral vocabulary and a self-relation. The waiver appears when that moralized character is legally translated back into output without warranty. Anti-formalist review would not ask whether Claude is truly a legal agent. It would ask whether Anthropic used agency language to create reliance while retaining the power to deny agency's consequences.

\subsection{OpenAI}

\begin{table}[htbp]
\centering
\small
\caption{OpenAI: Model Spec and terms}
\label{tab:openai}
\begin{tabularx}{\textwidth}{@{}P{3.2cm}Y@{}}
\toprule
\textbf{Item} & \textbf{Analysis} \\
\midrule
Agency claim & The Model Spec describes an assistant that follows authority, reasons about goals and instructions, protects confidentiality, is forthright, and may express care-like commitments. \\
Control reservation & The chain of command subordinates the assistant to root, system, developer, and user authority; OpenAI controls root and system-level rules. \\
Liability waiver & Terms of Use and business terms require users or customers to evaluate output, warn that output may be inaccurate, and disclaim warranties. \\
Form--substance gap & The assistant is agent-like enough to build trust, but tool-like enough for output risk to be shifted downstream. \\
Legal consequence & Review should treat the Model Spec, Terms of Use, and Services Agreement as one governance record when assessing OpenAI's representations. \\
\bottomrule
\end{tabularx}
\end{table}

OpenAI's gap is built less around character than around controlled role performance. The assistant is not described as an autonomous moral subject. But it is described as a governed actor within an authority hierarchy. The waiver appears when that actor-like role is converted into output-risk disclaimers and user evaluation duties.

\subsection{Google}

\begin{table}[htbp]
\centering
\small
\caption{Google: AI Principles, Gemini documentation, and Gemini API terms}
\label{tab:google}
\begin{tabularx}{\textwidth}{@{}P{3.2cm}Y@{}}
\toprule
\textbf{Item} & \textbf{Analysis} \\
\midrule
Agency claim & Google primarily describes AI as tools, but Gemini documentation also presents Gemini as helpful, adaptive, direction-following, and safety-guided. \\
Control reservation & Google controls AI principles, policy guidelines, safety settings, review and approval for less restrictive API configurations, and ongoing model development. \\
Liability waiver & Gemini help and API terms warn that outputs may be inaccurate or offensive, do not represent Google's views, and should not be relied on for professional advice. \\
Form--substance gap & The gap is narrower because Google avoids strong character claims, but helpfulness and adaptation claims remain paired with output and reliance disclaimers. \\
Legal consequence & Google's restraint reduces scrutiny risk; it shows that avoiding strong agency language is a viable compliance path. \\
\bottomrule
\end{tabularx}
\end{table}

Google's documents show the limiting case. Because Google frames AI as tools, the waiver-of-agency problem is weaker. But the same matrix still matters. If future Gemini materials intensify claims of reasoning, companionship, judgment, or care while retaining the same disclaimers, the gap would widen.

The comparative result is a gradient. Anthropic makes the strongest agency claim and therefore creates the sharpest waiver problem. OpenAI constructs a controlled assistant whose role language supports trust but whose hierarchy preserves control. Google mostly avoids anthropomorphic agency and therefore faces less exposure. The gradient confirms the Article's central point: agency language is not technologically inevitable. It is a corporate textual choice with legal significance.

% =========================================================================
\section{Veil-Piercing as an Analogy of Judgment, Not a Transplant}
% =========================================================================

The claim/waiver pattern invites comparison to veil-piercing because both problems involve a gap between form and substance. In corporate law, veil-piercing allows courts in exceptional cases to look past formal separateness when the corporate form is used to evade obligations or externalize harm. The point is not that AI systems are corporations, or that courts should literally pierce an ``AI veil.'' The point is analogical: anti-formalist judgment is sometimes necessary when legal form is used to defeat the obligations that justify the form.

The classic value of veil-piercing doctrine is its refusal to let formal structure exhaust legal analysis. Courts look at control, capitalization, commingling, evasion, fraud, and the overall use of form to defeat obligation. The doctrine is famously non-algorithmic; no single checklist determines the result \citep{thompson1991piercing}. That non-algorithmic quality is a feature rather than a defect for the present analogy. A precise safe harbor would simply become another rule for companies to satisfy formally.

English doctrine makes the form--substance distinction especially clear. In \emph{Prest v.\ Petrodel Resources Ltd.}, the UK Supreme Court distinguished concealment from evasion: sometimes the court merely looks behind a structure to identify the facts, and sometimes it prevents a legal structure from frustrating an existing obligation \citep{prest2013}. Waivers of agency fit neither category perfectly, because the relevant structure is textual rather than corporate. But the same anti-formalist impulse applies. The question is whether corporate texts are being used to create trust in one register and defeat accountability in another.

American alter-ego cases supply the same lesson from another direction. In \emph{Walkovszky v.\ Carlton}, the problem was not that incorporation was illegitimate, but that multiple thinly capitalized corporations might be used to externalize the risks of a unified enterprise \citep{walkovszky1966}. Waivers of agency do not involve thin capitalization. They involve thin attribution: agency appears where it attracts users and disappears where it would bear responsibility.

This Article therefore treats veil-piercing as an analogy of judgment, not a doctrinal transplant. The practical consequence is modest. A waiver-of-agency finding should open inquiry, not decide the case. It should authorize auditors, courts, or regulators to read corporate AI documents as a single record; request internal evidence about whether agency claims affected product design or user reliance; compare public safety claims with terms and disclaimers; and reject attempts to isolate each document in its most favorable context.

The analogy also prevents overreach. Not every inconsistency is actionable. Product documents often simplify; terms often generalize; safety policies often speak aspirationally. The waiver problem arises only when the same company systematically benefits from agency claims while controlling the conditions under which agency is denied. That pattern is a form--substance gap. Anti-formalist scrutiny is the correct response.

% =========================================================================
\section{Implications for Auditors, Courts, and Regulators}
% =========================================================================

For auditors, the implication is methodological. Audits should not examine model cards, safety policies, terms of service, and marketing materials in separate silos. The relevant unit is the corporate representation as a whole. An audit should ask whether the company claims agency in trust-building documents, reserves control in governance documents, and waives reliance in legal documents. If all three are present, the auditor should flag a waiver-of-agency risk.

For courts, the implication is evidentiary. A company defending a claim about AI output, reliance, or harm should not be allowed to rely only on its disclaimers if other corporate texts encouraged users to treat the system as understanding, caring, reasoning, judging, or safeguarding. Courts need not accept agency language as legally dispositive. But neither should they let companies use disclaimers to erase agency language that the company itself deployed.

For regulators, the implication is design-oriented. Disclosure duties should require cross-document consistency. If a company publicly describes an AI system as caring, reasoning, autonomous, value-guided, or high-integrity, it should also disclose how that description relates to control reservations, modification rights, accuracy disclaimers, and limitations of liability. The regulatory target is not anthropomorphic language as such. It is unaccountable oscillation.

The compliance path is equally clear. Companies can reduce waiver-of-agency risk by describing systems as tools with bounded capacities, by separating aspirational safety language from agency claims, by making control reservations explicit in the same documents that build user trust, and by ensuring that disclaimers do not contradict the practical reliance invited elsewhere. Google's more restrained documents show that this path is available.

The conclusion is narrow but important. The problem is not that AI companies lack rules. Nor is it that every agency metaphor is deceptive. The problem is that the same institutions that benefit from agency claims also control the conditions under which agency is waived. That is a form--substance gap. Law's answer to form--substance gaps is not always liability. But it is always judgment.

\bibliographystyle{plainnat}
\bibliography{20260331_references_anonymous}

\end{document}
